\scanpage{54}
\mspage{II 2}
% lines 1-6
Ce théorème sera démontré dans les sections suivantes, \deleted{par réduction au cas où $X$ est un ouvert d'une courbe algébrique projective et lisse sur $k$, auquel cas l'énoncé \unsure{résultera|résulte} \inserted{\unsure{du résultat classique déduit}} de l'hyp. de Riemann pour les courbes algébriques sur un corps fini, démontrée par A. Weil \unsure{que} dans ce cas.}

% lines 7-10
Nous allons dans le présent n° donner quelques conséquences. Notons que la suite exacte (1.1) implique une suite exacte
\[
(I^{\mathrm{ab}}_{\deleted{\unsure{\mathrm{ab}}}})_{\pi_0}\longrightarrow\pi^{\mathrm{ab}}_{\deleted{\unsure{\mathrm{ab}}}}\longrightarrow\pi_0^{\mathrm{ab}}\longrightarrow1
\]
\deleted{\unsure{et}} et en passant aux composantes $\ell$-primaires
\[
I(\ell)_{\pi_0}\longrightarrow\pi^{\mathrm{ab}}(\ell)\mathrel{\text{---}}\pi_0^{\mathrm{ab}}(\ell)\mathrel{\text{---}}1,
\]
d'où on déduit :

% lines 14-17
\uline{Corollaire} 1.2. L'homomorphisme \inserted{surjectif} canonique
\[
\pi^{\mathrm{ab}}(\ell)\longrightarrow\pi_0^{\mathrm{ab}}(\ell)
\]
\deleted{\unsure{est}} a un \uline{noyau fini}.

\deleted{Utilisant I 6.1, on en déduit de 1.1 le}

% lines 19-27
\uline{Corollaire} 1.3. Soit
\[
u:\pi_1(X,\unsure{\xi|\zeta})\longrightarrow\operatorname{Aut}_{\mathbb{Q}_\ell}(\unsure{E|\mathbb{E}}\unsure{]|)}
\]
une représentation $\ell$-adique \deleted{\unsure{de $\pi$}}, tel que la représentation induite sur \deleted{\unsure{$I$}} \inserted{$\pi_1(\overline X,\xi)$} soit \inserted{\uline{essentiellement résoluble}} \deleted{\uline{abélienne}, \unsure{ou plus} généralement \deleted{\uline{commutative}} \uline{\unsure{abélienne}}} (i.e. $\exists$ un sous-groupe ouvert de $I$ dont l'action soit \deleted{abélienne} \inserted{résoluble}). Alors la représentation induite sur $I$ est \uline{essentiellement unipotente}, \deleted{et que} (donc essentiellement triviale si elle est semi-simple).
% archive sheet number: 53
